\section{A boundary that remembers transport}\label{ntg:sec:boundary}

We return to the fibre lattice $\Lambda$ in Section~1 and write
$L=\Lambda$ in its ordered marking $B_0=(e_1,e_2,e_3,e_4)$. Let
$\T^4=(L\otimes\mathbb R)/L$. An invertible integer matrix $T$ acts
on the torus by $[x]\mapsto[Tx]$. Its mapping torus is
\[
 Y_T=(\T^4\times[0,1])/((x,1)\sim(Tx,0)).
\]
We orient the interval in the increasing direction and put the interval
before the orientation of a transported fibre chain. A prism over a
cycle $z$ then has boundary $Tz-z$. Thus a nontrivial $T$ changes the
relations among higher-dimensional cycles, even when a chosen circle
inside the fibre is fixed.

The three monodromy matrices from Section~1, now all in the common
marking $B_0$, act on column vectors in $L$:
\begin{equation}\label{ntg:eq:native}
\begin{aligned}
 T_0&=\begin{pmatrix}0&-1&-1&0\\1&-1&0&0\\0&0&1&0\\0&0&0&1\end{pmatrix},
 &T_1&=\begin{pmatrix}-1&1&-1&2\\-1&0&-1&1\\1&1&2&-1\\0&0&0&1\end{pmatrix},\\
 T_\infty&=\begin{pmatrix}2&1&1&0\\0&1&0&-1\\-1&-1&0&1\\0&0&0&1\end{pmatrix}.
\end{aligned}
\end{equation}
They have determinant one, and direct multiplication gives
\[
 T_0^3=T_1^4=I,\qquad T_0T_1T_\infty=I,
 \qquad (T_\infty-I)^2=0.
\]
These are legitimate transport maps for a flat torus bundle over an
oriented sphere with three disks removed. Choose paths from one base
point to its boundary circles. Orient and order the resulting based
loops so that their relation is $\gamma_0\gamma_1\gamma_\infty=1$.
Our word convention is that the product $uv$ acts by $M_uM_v$;
equivalently, geometric paths in this notation are traversed from
right to left. Reversing a boundary parametrization inverts its
transport and reverses its interval orientation.

The difference from a product is visible in the third matrix alone.
Put
\begin{equation}\label{ntg:eq:shear-basis}
 a=e_1-e_3,\qquad b=e_1,\qquad c=e_3-e_2,\qquad d=e_4.
\end{equation}
The matrix with these columns has determinant one. This is a different
basis from the determinant-minus-one basis (1.11) for the alternating
form. In the present integral basis the action is
\begin{equation}\label{ntg:eq:shear}
 T_\infty a=a,\quad T_\infty b=b+a,\quad
 T_\infty c=c,\quad T_\infty d=d+c.
\end{equation}
We will prove
\[
 H_2(Y_{T_\infty};\Z)=H_3(Y_{T_\infty};\Z)=\Z^6.
\]
For the identity matrix the mapping torus is $\T^5$, and both ranks
are ten. Any filling calculation using a product in place of this
boundary therefore requires a comparison map; agreement on one fixed
circle cannot provide it.

The full gluing problem has one further ingredient. An inclusion of a
boundary mapping torus into a filling contains information about both
fibre chains and their transported prisms. We first describe exactly
how this information enters the differential, then compute the boundary
groups and their torsion.

\section{Chains over a graph}

A chain complex of free abelian groups is a collection $C_j$, zero
for $j<0$, with homomorphisms $d:C_j\to C_{j-1}$ satisfying $d^2=0$.
Its homology is $H_j(C)=\ker d_j/\im d_{j+1}$. A chain map commutes
with $d$. A degree-one map $H:C_j\to Q_{j+1}$ is a chain homotopy
from $g$ to $f$ when $d_QH+Hd=f-g$.

For a space we may use singular chains, the free groups generated by
continuous maps from oriented simplices, with the alternating face
boundary. A homeomorphism acts on these chains strictly: composition
of homeomorphisms is composition of chain maps. This avoids making
simultaneous choices of cellular approximations for several transport
maps.

Here $C$ denotes a chain complex; the change-of-marking matrix in
(1.5) retains its meaning only in the earlier lattice calculation.
Fix chain automorphisms $M_0,M_1$ of $C$, and set
$M_\infty=(M_0M_1)^{-1}$. They are the chain form of transport.
On the fibre homology $H_k(\T^r;\Z)=\bigwedge^k\Z^r$, transport by
an integer matrix $T_i$ induces $\bigwedge^kT_i$; a proof of this
identification appears in Section~\ref{ntg:sec:torus}. On cohomology the
contragredient transport is $T_i^{-\mathsf T}$ in degree one. The
homological and cohomological conventions should not be interchanged.

\begin{definition}\label{ntg:def:prism}
The chain complex of one loop with transport $M_i$ is
\[
 (B_i)_n=C_n\oplus C_{n-1},\qquad
 d_{B_i}(u,v)=(du+(M_i-I)v,-dv).
\]
The first component is in a fibre over the vertex; the second is the
oriented prism over a fibre chain. For the graph consisting of two
loops with a common vertex, put
\begin{equation}\label{ntg:eq:graph}
 E_n=C_n\oplus C_{n-1}\oplus C_{n-1},\qquad
 d_E(a,b_0,b_1)=(da+(M_0-I)b_0+(M_1-I)b_1,-db_0,-db_1).
\end{equation}
\end{definition}

The signs follow from
$\partial([0,1]\times z)=(1\times z)-(0\times z)-[0,1]\times\partial z$.
After identifying the ends of the prism, the first two terms become
$(M_i-I)z$. Since $M_i$ commutes with $d$, both displayed
differentials square to zero.

\begin{lemma}\label{ntg:lem:boundary-word}
The three boundary maps into the graph complex are the chain maps
\begin{equation}\label{ntg:eq:boundary-word}
 \begin{split}
 J_0(u,v)&=(u,v,0),\\
 J_1(u,v)&=(u,0,v),\\
 J_\infty(u,v)&=(u,-M_0^{-1}v,-M_\infty v).
 \end{split}
\end{equation}
The last map represents the boundary word
$\gamma_\infty=(\gamma_0\gamma_1)^{-1}$ with the preceding
orientation convention.
\end{lemma}
\begin{proof}
An inverse path transports a chain by $M_i^{-1}$ and traverses its
prism backwards. Its coefficient in the original oriented edge is
therefore $-M_i^{-1}$. Traversing $\gamma_0^{-1}$ followed by
$\gamma_1^{-1}$ gives coefficient $-M_0^{-1}$ in the first edge and
$-M_1^{-1}M_0^{-1}=-M_\infty$ in the second. The endpoints cancel
at the intermediate vertex. Algebraically this is the identity
\begin{equation}\label{ntg:eq:word-identity}
 (M_0-I)(-M_0^{-1})+(M_1-I)(-M_\infty)=M_\infty-I,
\end{equation}
because $M_1M_\infty=M_0^{-1}$. This identity verifies the vertex
component of $d_EJ_\infty=J_\infty d_{B_\infty}$. The two edge
components agree because $M_0^{-1}$ and $M_\infty$ commute with $d$.
The other two chain-map identities follow directly from their formulas.
\end{proof}

We explain the topological content of these complexes. Cut a loop at
its vertex. The bundle becomes a product over an interval; its two end
inclusions are the identity and transport. Chains on the vertex fibre,
together with prisms, give the complex $B_i$. For two loops the same
construction gives $E$. It applies to a sphere with three boundary
circles after a deformation retraction to its two-loop graph, retaining
the chosen paths to the boundary circles.

One can verify this chain model without assuming a special subdivision
of the torus. Use collars of the cut fibres. Subdivide singular simplices
until each lies in a collar or in an interval product. The alternating
sum of the subdivision prisms is a chain homotopy between subdivision
and the identity. The interval products retract to a fibre; keeping the
retraction prisms leaves one copy of $C$ at each vertex and one shifted
copy along each edge. Their boundary is exactly the end difference
minus the prism over $d$. Gluing the cuts gives
\eqref{ntg:eq:graph}. For a boundary path, the successive interval prisms
give the coefficients in \eqref{ntg:eq:boundary-word}. A lift to the
universal covering tree makes the cancellation literal: each traversed
edge occurs once with its transported coefficient. This also gives a
chain comparison for the boundary maps, not just an isomorphism of
the resulting homology groups.

Changing a retraction or a subdivision changes this comparison by a
chain homotopy. The filling maps must be transported through the same
comparison. This compatibility will be part of the geometric
hypotheses of the gluing theorem.

\section{What a filling must specify}

Let $Q_i$ be a free chain complex for a proposed filling of the $i$th
boundary, for $i=0,1,\infty$. We write $d_i$ for its differential.
The map from the boundary complex to $Q_i$ consists of maps
\begin{equation}\label{ntg:eq:filling-data}
 F_i:C_n\longrightarrow(Q_i)_n,\qquad
 H_i:C_n\longrightarrow(Q_i)_{n+1}.
\end{equation}
Here $F_i$ is inclusion of a fibre, and $H_i$ describes how its
transport prism lies in the filling. Define
$K_i(u,v)=F_iu+H_iv$. The condition that $K_i$ commute with
boundaries is precisely
\begin{equation}\label{ntg:eq:homotopy}
 d_iF_i=F_id,\qquad d_iH_i+H_id=F_i(M_i-I).
\end{equation}
Indeed, comparison of the coefficients of $u$ and $v$ in
$d_iK_i(u,v)=K_id_{B_i}(u,v)$ gives these two equations separately.

In particular, $F_i$ kills monodromy minus the identity on homology.
The dual pullback from the filling therefore lands in invariant fibre
cohomology. Specifying the image of that pullback is weaker than
specifying $F_i$, and even the chain map $F_i$ does not specify $H_i$.
The second equation of \eqref{ntg:eq:homotopy} is the next attachment
equation that a filling has to solve.

\begin{theorem}[The gluing differential]\label{ntg:thm:gluing}
Let $C,Q_i,M_i,F_i,H_i$ satisfy the hypotheses and equations above.
Set
\[
 A=E\oplus\bigoplus_iQ_i,\qquad B=\bigoplus_i B_i,
 \qquad L((w_i)_i)=\bigl(\textstyle\sum_iJ_iw_i,(-K_iw_i)_i\bigr).
\]
Then $L:B\to A$ is a chain map, and the groups
$G_n=A_n\oplus B_{n-1}$ with differential
\begin{equation}\label{ntg:eq:cone}
 D(x,y)=(d_Ax+Ly,-d_By)
\end{equation}
form a chain complex. In coordinates
\[
 (a,b_0,b_1,(q_i)_i,(u_i)_i,(v_i)_i)\in G_n,
\]
where $a\in C_n$, $b_j,u_i\in C_{n-1}$,
$q_i\in(Q_i)_n$, and $v_i\in C_{n-2}$, its differential is
\begin{equation}\label{ntg:eq:full-boundary}
\begin{aligned}
 a'&=da+(M_0-I)b_0+(M_1-I)b_1+u_0+u_1+u_\infty,\\
 b'_0&=-db_0+v_0-M_0^{-1}v_\infty,\\
 b'_1&=-db_1+v_1-M_\infty v_\infty,\\
 q'_i&=d_iq_i-F_iu_i-H_iv_i,\\
 u'_i&=-du_i-(M_i-I)v_i,\\
 v'_i&=dv_i.
\end{aligned}
\end{equation}
If these chain data represent a torus bundle over the three-holed
sphere, its three boundary inclusions, and inclusions into three
fillings, then $H_*(G)$ is the integral homology of their gluing by
collars. In particular this conclusion applies to ordinary gluing
along collared boundaries. For arbitrary attachment maps it computes
the gluing with mapping cylinders inserted.
\end{theorem}
\begin{proof}
Lemma~\ref{ntg:lem:boundary-word} and \eqref{ntg:eq:homotopy} show that
$d_AL=Ld_B$. Consequently
\[
 D^2(x,y)=(d_A^2x+d_ALy-Ld_By,d_B^2y)=0.
\]
Substitution gives \eqref{ntg:eq:full-boundary}. It is useful to check the
two components where the attachment data matter. In the vertex
component, the remaining coefficient of $v_\infty$ is
\[
 -(M_0-I)M_0^{-1}-(M_1-I)M_\infty-(M_\infty-I)=0.
\]
In the $i$th filling component it is
$-d_iH_i-H_id+F_i(M_i-I)=0$. The other terms cancel by
$d^2=d_i^2=0$ and commutation of transport with $d$.

For the geometric assertion, replace each attachment by its mapping
cylinder. A cylinder has two end copies of the boundary chains and a
prism over them. The end copies attach by $J_i$ and $K_i$; its oriented
boundary gives $J_i-K_i$. The boundary on the prism summand is $-d_B$.
After retracting the two halves of the collars to their ends, these are
exactly the summands and differential of \eqref{ntg:eq:cone}.
The same subdivision and prism argument used for the graph applies
to this two-piece decomposition: sufficiently small chains belong to
one side or the overlap, and the duplicate overlap chains are removed
by their cylinder prisms. It produces a chain comparison inducing
isomorphisms on homology. If the attaching maps are boundary inclusions
with collars, compressing the inserted cylinders gives a homotopy
equivalence to the usual glued space. For other maps, the cylinders
are part of the stated construction.
\end{proof}

For clarity, in degree two the variables are
\[
 a\in C_2,\quad b_j,u_i\in C_1,\quad q_i\in(Q_i)_2,
 \quad v_i\in C_0.
\]
Thus $H_i:C_0\to(Q_i)_1$ already occurs in $D_2$. In degree three,
$v_i\in C_1$, so $H_i:C_1\to(Q_i)_2$ occurs in $D_3$.
Computing $H_3$ also requires $D_4$, hence
$H_i:C_2\to(Q_i)_3$. The formula neither discards these terms nor
replaces them by their induced maps on homology.

\begin{proposition}[Compatible maps]\label{ntg:prop:natural}
Suppose $f:C\to\widetilde C$ and $g_i:Q_i\to\widetilde Q_i$ are
chain maps with
\[
 fM_i=\widetilde M_if,\qquad
 g_iF_i=\widetilde F_if,\qquad
 g_iH_i=\widetilde H_if.
\]
Applying $f$ to all fibre and prism coordinates and $g_i$ to $q_i$
gives a chain map $G\to\widetilde G$. These maps preserve identities
and composition.
\end{proposition}
\begin{proof}
The first identity also holds for inverses and products of transport
maps. Each of the six lines of \eqref{ntg:eq:full-boundary} therefore
commutes with the proposed map. Identity maps and composites satisfy
the same three equations.
\end{proof}

There is also a controlled change of attachment representative. If
$L'-L=d_AS+Sd_B$ for a degree-one map $S:B\to A$, then
\[
 \Cone(L')\longrightarrow\Cone(L),\qquad (x,y)\longmapsto(x+Sy,y)
\]
is a chain isomorphism with inverse $(x,y)\mapsto(x-Sy,y)$.
Substituting into \eqref{ntg:eq:cone} proves the assertion. Thus a
specified homotopy between attachments permits a comparison. Equality
of maps on homology alone does not supply such a homotopy.

\begin{example}[The prism map can change torsion]\label{ntg:ex:missing-H}
Let $C_0=\Z$, $C_j=0$ for $j\ne0$, and $M=I$. Then $B$ is the
cellular complex of a circle, with one vertex and one edge.
Take $Q_0=Q_1=\Z$ and $d_Q=0$, also a circle. Set $F_0=1$ and
let $H_0$ be multiplication by an integer $k$.
Equation~\eqref{ntg:eq:homotopy} holds for every $k$. The resulting
$K:B\to Q$ is the chain map of a degree-$k$ map of circles.
Its mapping cone has differential $k$ from the degree-two prism to
the degree-one circle, while the degree-one prism over the vertex
maps isomorphically to the degree-zero vertex. Hence
\[
 H_1(\Cone K)=\Z/k\Z,\qquad
 H_2(\Cone K)=\ker(k:\Z\to\Z).
\]
These are the reduced homology groups of the circle with a disk
attached by the degree-$k$ map. The map $F$ is identical for all $k$.
For $k=0$ both displayed groups are $\Z$; for $|k|=1$ both vanish.
The missing prism map changes the homology, even with the fibre
inclusion fixed at chain level.
\end{example}

\section{Integral torsion and prime coefficients}\label{ntg:sec:torsion}

When $C$ and the $Q_i$ are finite free in each degree, the differential
in \eqref{ntg:eq:full-boundary} is an integer matrix in every degree.
There is a direct way to read its integral information without changing
coefficients prematurely.

\begin{lemma}\label{ntg:lem:smith}
An integer matrix can be changed, by invertible integer row and column
operations, to a diagonal matrix with positive nonzero entries
$s_1,\ldots,s_t$ satisfying $s_1\mid s_2\mid\cdots\mid s_t$,
and all other entries zero. The gcd of all its $j$-rowed minors is
$s_1\cdots s_j$ for $j\le t$; all larger minors vanish.
\end{lemma}
\begin{proof}
Move a nonzero entry to the first diagonal position. Divide entries
in its column by it, subtract multiples of the first row, and interchange
rows when a smaller nonzero remainder occurs. Do the corresponding
operations on its row. Positive remainders strictly decrease, so this
process reaches an entry that clears its row and column. If that
entry does not divide some entry in the remaining block, adding that
entry's row to the first row and repeating Euclidean division produces
a smaller positive pivot. This also terminates. The final pivot
divides the remaining block. Induction on the block gives the diagonal
form and the divisibility chain. These operations and their inverses
have integer entries and determinant $\pm1$.

Each minor after an elementary operation is an integer linear
combination of the old minors of the same size. Applying the inverse
operation proves equality of their gcds. For the diagonal form the gcd
is the product of its first $j$ entries, by divisibility. This proves
the formula.
\end{proof}

To compute $H_n(G)$, first apply the lemma to $D_n$ and retain the
integer change of basis in $G_n$. The columns corresponding to zero
diagonal entries give an integral basis of $\ker D_n$. The equation
$D_nD_{n+1}=0$ puts every column of $D_{n+1}$ in that basis. Apply
the lemma again to this restricted matrix, say with nonzero factors
$s_1,\ldots,s_t$. If $z=\rank\ker D_n$, then
\begin{equation}\label{ntg:eq:homology-algorithm}
 H_n(G)=\Z^{z-t}\oplus\bigoplus_{j=1}^t\Z/s_j\Z.
\end{equation}
This is a finite calculation for specified finite chain data. A single
Smith form of $D_{n+1}$ in an unrelated basis is insufficient; its
image must be expressed inside the integral kernel of $D_n$.

Let $p$ be prime. Reduction modulo $p$ gives a complex $G/pG$.
For a mod-$p$ cycle $\bar x$ in degree $n$, choose an integral lift
$x$. There is an integral chain $y$ with $Dx=py$, since every
coordinate of $Dx$ is divisible by $p$. Define
\begin{equation}\label{ntg:eq:bockstein}
 \beta_p:H_n(G/pG)\longrightarrow H_{n-1}(G),\qquad
 \beta_p[\bar x]=[Dx/p].
\end{equation}
Freeness of the chain groups gives $Dy=0$ from $pDy=D^2x=0$.
Replacing $x$ by $x+pz$ changes $y$ by $Dz$. Replacing $\bar x$
by a homologous cycle has the same effect after lifting its boundary.
Thus \eqref{ntg:eq:bockstein} is well-defined. Its image consists exactly
of the classes killed by $p$: its values are killed by $p$, and if
$p[y]=0$, choose $x$ with $Dx=py$ and reduce $x$ modulo $p$.
This is the integral connecting homomorphism, often called the
Bockstein homomorphism.

In a two-term direct summand $\Z x\xrightarrow{s}\Z y$, a prime
$p\mid s$ gives a mod-$p$ class $[\bar x]$ and
\begin{equation}\label{ntg:eq:beta-factor}
 \beta_p[\bar x]=(s/p)[y]\in\Z/s\Z.
\end{equation}
For $s\ne0$ this element has order $p$. If $p$ does not divide $s$,
the two-term complex is acyclic modulo $p$. If $s=0$, both classes
are free integrally and their connecting homomorphisms vanish. These
three cases distinguish torsion, an invertible relation, and a genuine
rank increase.

\section{The torus mapping-torus complex}\label{ntg:sec:torus}

The torus $\T^r$ has one vertex and one oriented edge in each circle
factor. Taking products gives one $k$-cell for each increasing
$k$-tuple of coordinate circles. Opposite faces cancel, so the
cellular differential is zero. Denote the resulting group of cells by
$L_k=\bigwedge^k\Z^r$, with basis
$e_{i_1}\wedge\cdots\wedge e_{i_k}$. It follows directly that
$H_k(\T^r;\Z)=L_k$.

The map $[x]\mapsto[Tx]$ induces $\bigwedge^kT$ on these groups.
Indeed, its coefficient from the coordinate torus indexed by columns
$J$ to the class indexed by rows $I$ is found by projecting onto the
coordinates in $I$. The resulting map between oriented $k$-tori has
integer matrix $T_{I,J}$ and degree $\det T_{I,J}$. For a nonsingular
integer matrix, this follows by counting its kernel, the quotient of
the integer lattice by its image, with the common local orientation
sign; integer row and column operations reduce the count to a
diagonal matrix. A singular matrix factors through a real subspace of
smaller dimension, so its degree is zero. These are exactly the
minors defining the exterior power.

\begin{proposition}\label{ntg:prop:mappingtorus}
For $T\in\mathrm{GL}(r,\Z)$, the additive integral homology of
$Y_T$ is computed by the finite free chain complex
\begin{equation}\label{ntg:eq:torus-cone}
 P_n=L_n\oplus L_{n-1},\qquad
 d_P(x,y)=((\bigwedge^{n-1}T-I)y,0).
\end{equation}
Consequently
\begin{equation}\label{ntg:eq:mapping-homology}
 H_n(Y_T;\Z)\cong
 \coker(\bigwedge^nT-I)\oplus
 \ker(\bigwedge^{n-1}T-I),
\end{equation}
where exterior powers outside $0,\ldots,r$ are zero.
\end{proposition}
\begin{proof}
Use the product cell structure just described and replace the single
torus map by a homotopic cellular map. Such a replacement is made
one skeleton at a time: move the image of each cell off the interiors
of target cells of larger dimension by a small deformation and radial
retraction, extending the deformation over the cell using its collar.
The deformations agree on faces by induction. A homotopy of the map
gives a homotopy equivalence of the two mapping-torus constructions,
by inserting the homotopy as a cylinder between their end gluings.

The cellular map acts by $\bigwedge^kT$, because cellular chains
already equal cellular homology when all differentials vanish.
Attach one prism for every fibre cell. Its ends contribute
$(\bigwedge^kT-I)y$ and its side boundary is zero in cellular
chains. This proves \eqref{ntg:eq:torus-cone}. A cycle consists of an
arbitrary $x\in L_n$ and $y\in\ker(\bigwedge^{n-1}T-I)$.
Boundaries lie entirely in the first component and form
$\im(\bigwedge^nT-I)$. Taking the quotient gives
\eqref{ntg:eq:mapping-homology}.
\end{proof}

This proposition concerns one transport map and additive homology.
It does not identify singular chains simultaneously with
$\bigwedge^*L$ while preserving all transport maps, filling maps,
and their chain homotopies. The general gluing theorem retains $C$
and its differential. A finite exterior-power replacement in that
theorem requires compatible comparisons for every displayed map.

\section{Two integral shears and their covers}\label{ntg:sec:shears}

For $h\in\Z$, let $T_h=T_\infty^h$. Since
$(T_\infty-I)^2=0$, we have $T_h=I+h(T_\infty-I)$, also for
negative $h$. The basis \eqref{ntg:eq:shear-basis} gives
\[
 T_ha=a,\qquad T_hb=b+ha,\qquad
 T_hc=c,\qquad T_hd=d+hc.
\]
For $h>0$, $Y_{T_h}$ is the pullback of $Y_{T_\infty}$ through
the degree-$h$ cover of its base circle. Concatenating $h$ interval
trivializations multiplies their end maps, which proves the claim.
For $h<0$ one also reverses the parametrization of the base circle.
For $h=0$ we simply mean the identity transport and $\T^5$.

Write $ab$ for $a\wedge b$, and use the same convention for longer
products. Set $A_k=\bigwedge^kT_h-I$. In degree one,
\begin{equation}\label{ntg:eq:degreeone}
 A_1b=ha,\qquad A_1d=hc,\qquad A_1a=A_1c=0.
\end{equation}
In degree two the complete list of nonzero images is
\begin{equation}\label{ntg:eq:degreetwo}
 A_2(ad)=h\,ac,\qquad A_2(bc)=h\,ac,
 \qquad A_2(bd)=h(ad+bc)+h^2ac.
\end{equation}
Indeed, $(b+ha)\wedge c=bc+hac$, while
$(b+ha)\wedge(d+hc)=bd+hbc+had+h^2ac$.
The other three basis elements $ab,ac,cd$ are fixed. The quadratic
term comes from transporting both factors of $bd$; discarding it
would change the chain map, even though it will not change its Smith
factors. In degree three,
\begin{equation}\label{ntg:eq:degreethree}
 A_3(abd)=h\,abc,\qquad A_3(bcd)=h\,acd,
 \qquad A_3(abc)=A_3(acd)=0.
\end{equation}
Finally, $A_0=A_4=0$.

\begin{lemma}\label{ntg:lem:shear-smith}
For $h\ne0$, each of $A_1,A_2,A_3$ has precisely two nonzero
Smith factors, both equal to $|h|$. Their kernel ranks are respectively
$2,4,2$. The images are
\[
 h\langle a,c\rangle,\qquad
 h\langle ac,ad+bc\rangle,\qquad
 h\langle abc,acd\rangle.
\]
Each pair in angle brackets extends to an integral basis of the
corresponding exterior lattice; such a sublattice is called primitive.
\end{lemma}
\begin{proof}
Equations~\eqref{ntg:eq:degreeone} and \eqref{ntg:eq:degreethree} are already
diagonal after ordering their source and target bases. In degree two,
replace the source vector $bd$ by $bd-had$. Its image is
$h(ad+bc)$. The six source vectors
\[
 ad,\quad bd-had,\quad ab,\quad ac,\quad cd,\quad ad-bc
\]
form an integral basis: they recover $bc=ad-(ad-bc)$ and
$bd=(bd-had)+had$, as well as the other four original basis vectors.
The last four vectors lie in the kernel. In the target, $ac$ and
$ad+bc$ extend to an integral basis by adding $ab,ad,bd,cd$.
In these independent source and target bases the only nonzero entries
are $h,h$. Changing signs makes them positive. This proves all
assertions, including primitivity. Notice that the source and target
bases are different; no decomposition into invariant and image
lattices over $\Z$ has been assumed.
\end{proof}

\begin{theorem}\label{ntg:thm:shear-homology}
For $h\ne0$, put $A_h=\Z/|h|\Z$. The integral homology groups
of $Y_{T_h}$ are
\begin{equation}\label{ntg:eq:shear-homology}
\begin{array}{c|cccccc}
 n&0&1&2&3&4&5\\ \hline
 H_n(Y_{T_h};\Z)&\Z&\Z^3\oplus A_h^2&\Z^6\oplus A_h^2&
 \Z^6\oplus A_h^2&\Z^3&\Z.
\end{array}
\end{equation}
All other homology groups vanish. For $h=0$, they are
$\Z^{\binom5n}$ in degrees $0\le n\le5$.
\end{theorem}
\begin{proof}
The ranks of $L_k$ are $1,4,6,4,1$. For $h\ne0$, the ranks of
$A_k$ are $0,2,2,2,0$. Subtracting gives cokernel free ranks
$1,2,4,2,1$ and kernel ranks $1,2,4,2,1$.
The only torsion is the pair of factors in
Lemma~\ref{ntg:lem:shear-smith}, in cokernel degrees $1,2,3$.
Adding the cokernel in degree $n$ and the kernel in degree $n-1$
as in \eqref{ntg:eq:mapping-homology} gives the table. When $h=0$ every
differential in \eqref{ntg:eq:torus-cone} vanishes. Its degree-$n$ rank
is $\binom4n+\binom4{n-1}=\binom5n$.
\end{proof}

At $h=1$ the integral homology is torsion-free, with free ranks, or
Betti numbers, $(1,3,6,6,3,1)$. Thus the transport in
\eqref{ntg:eq:shear} can remove four
free classes from each of degrees two and three without introducing
torsion. At $h=2$ the same free ranks remain, but there are two
order-two classes in each of degrees one, two, and three. At $h=6$
these become two order-six classes. The index of the image lattice,
rather than its rational span, distinguishes the three cases.

\begin{proposition}\label{ntg:prop:shear-beta}
Let $h\ne0$ and let $p$ be prime. If $p\nmid h$, the dimensions
of $H_n(Y_{T_h};\F_p)$ are $(1,3,6,6,3,1)$ and $\beta_p=0$.
If $p\mid h$, they are $(1,5,10,10,5,1)$.
In this second case $\beta_p$ has rank two over $\F_p$ from
degree $j+1$ to the order-$p$ subgroup of $H_j(Y_{T_h};\Z)$,
for each $j=1,2,3$, and is zero in the other degrees.
\end{proposition}
\begin{proof}
Use the diagonal bases of Lemma~\ref{ntg:lem:shear-smith} in the fibre
and suspended summands of \eqref{ntg:eq:torus-cone}. They are independent
bases of independent chain summands, so they decompose the complex
into free zero-differential summands and six two-term complexes with
differential $h$. For $p\nmid h$ these six summands are acyclic
modulo $p$. For $p\mid h$ the entire displayed differential is zero
modulo $p$, giving the ranks of the five-torus chain groups.
Formula~\eqref{ntg:eq:beta-factor} proves the rank assertion. More
explicitly, write $\sigma z=(0,z)$ for a suspended fibre chain.
The two source and target pairs in each degree are
\begin{equation}\label{ntg:eq:explicit-beta}
\begin{array}{c|c|c}
 j&\text{mod-$p$ sources in degree }j+1&
       \text{integral targets in degree }j\\ \hline
 1&\sigma b,\ \sigma d&a,\ c\\
 2&\sigma(ad),\ \sigma(bd-had)&ac,\ ad+bc\\
 3&\sigma(abd),\ \sigma(bcd)&abc,\ acd.
\end{array}
\end{equation}
Each listed source has integral boundary $h$ times its listed
target, so its connecting image is $(h/p)$ times that target.
These are independent classes of order $p$. All remaining
mod-$p$ classes have integral cycle lifts in the diagonal complex
and hence connecting image zero.
\end{proof}

At $h=p=2$, the degree-three suspended chain has boundary
\[
 \partial\sigma(bd-2ad)=2(ad+bc),\qquad
 \beta_2[\overline{\sigma(bd-2ad)}]=[ad+bc].
\]
The coefficient of $bc$ modulo two defines a homomorphism
$L_2\to\F_2$ that vanishes on $\im A_2$ and takes $ad+bc$ to one.
Thus $[ad+bc]$ is nonzero in $\coker A_2$ and has order exactly two.
Define $\ell:L_2\to\Z$ by
$\ell(ad)=1$, $\ell(bc)=-1$, and zero on $ab,ac,bd,cd$.
It vanishes on $\im A_2$ and hence descends to the cokernel.
Since $\ell(ad-bc)=2$, the class $[ad-bc]$ has infinite order.

The degree-two class $[\overline{\sigma b}]$ at
$h=6$ maps under $\beta_2$ to $3[a]\in\Z/6\Z$ and under
$\beta_3$ to $2[a]$. The degree-three classes over $ad$ and
$bd-6ad$ detect the two torsion generators in $H_2$ in exactly
the same way. These formulas retain the integral normalization of
the cycles, including the quadratic term in \eqref{ntg:eq:degreetwo}.

\section{The remaining attachment equations}\label{ntg:sec:remaining}

In the geometric setting of Section~5, choose disjoint closed disks
$D_i\subset\mathbb P^1$ around $i=0,1,\infty$ and set
\[
 \Sigma=\mathbb P^1\setminus\bigcup_i\operatorname{int}D_i,
 \qquad N_i=f^{-1}(D_i),\qquad E_f=f^{-1}(\Sigma).
\]
Choose collars of the boundary circles and paths from a common base
point, with the loop and interval orientations of
Section~\ref{ntg:sec:boundary}. The five-dimensional space
$Y_{T_i}$ denotes the corresponding boundary mapping torus. It is
distinct from the six-dimensional total space $Y(u)$.
An application of Theorem~\ref{ntg:thm:gluing} to $Y(u)$ requires
chain comparisons for $E_f$, each $\partial N_i$, and each $N_i$,
compatible with these boundary identifications.

For each $i$, $Q_i$ must represent integral chains on $N_i$, with its
actual differential $d_i$. The map $F_i:C_n\to(Q_i)_n$ must represent
the inclusion of a nearby fibre into $N_i$. The map
$H_i:C_n\to(Q_i)_{n+1}$ must represent its oriented boundary prism
inside $N_i$. In particular, $F_i$ and $H_i$ are required with their
integer coefficients and chain comparisons, not only with their induced
maps on cohomology. The necessary degrees are $0\le n\le4$ for
$Q_i$, $0\le n\le3$ for $F_i$, and $0\le n\le2$ for $H_i$,
together with the fibre complex through degree four and the transport
maps in those degrees. These data determine $D_2,D_3,D_4$.

At $0$ and $1$, the logarithmic classes $a_0,a_1$ in (5.1)
and the affine actions $z\mapsto T_iz+a_i$ are part of the attachment
data. Proposition~\ref{prop:affine-free} gives their exact freeness
conditions. A finite chain model for the quotient neighbourhood must
retain these translations as well as the linear transport. At infinity,
the two-shear calculation determines the boundary; a chain model for
the degeneration and its nearby-fibre inclusion is still required.
The first-degree transgression (5.4) and the fundamental-group
presentation (5.6) retain the geometric hypotheses and external
dependencies stated in Section~5. Neither calculation specifies these
higher chain maps.

The matrices \eqref{ntg:eq:native} determine the flat torus bundle over
the three-holed sphere and all three boundary mapping-torus complexes.
They do not determine the spaces that fill those boundaries.
To complete a particular filling calculation one must specify, for
each $i$, the chain groups $(Q_i)_n$, their boundary maps, and
the maps $F_{i,n},H_{i,n}$ in \eqref{ntg:eq:filling-data}, with the
following equations in every required degree:
\begin{equation}\label{ntg:eq:missing-equations}
 \begin{split}
 d_{i,n}F_{i,n}&=F_{i,n-1}d_n,\\
 d_{i,n+1}H_{i,n}+H_{i,n-1}d_n
     &=F_{i,n}(M_{i,n}-I).
 \end{split}
\end{equation}
Here $M_{i,n}$ means the chain action in degree $n$, not merely
its induced action on homology. For $H_2$ and $H_3$ the resulting
matrices $D_2,D_3,D_4$ and their common bases are needed. Their
kernel and image calculation is \eqref{ntg:eq:homology-algorithm}.

A degeneration with a toric central fibre offers a concrete way to
construct $Q_\infty$: choose chains on the components, on their
pairwise intersections, and on their multiple intersections, and
record every inclusion and every gluing identification with signs.
The alternating face maps and the internal boundaries form a total
complex. If a component is identified with itself along different
boundary curves, those two maps must be retained separately. A
description of the incidence pattern alone gives neither their
integer matrices nor the specialization prism from a nearby smooth
torus. The latter is exactly the information in
$F_\infty,H_\infty$, together with a comparison of this total
complex to the filling. The equations in
\eqref{ntg:eq:missing-equations} express the required compatibility.

At a finite-order boundary, the linear part of an affine torus action
also does not determine its quotient. For instance, translation by
$a$ on a circle has order dividing $m$ when $ma\in\Z$; its
order is smaller if the denominator of $a$ is smaller. In a higher
torus an affine map $x\mapsto Tx+a$ has a fixed point precisely
when $(T-I)x+a$ lies in the integral lattice. One must check every
nonidentity group element to conclude that a quotient is free.
If the action is not free, a singular quotient, its resolution, and
a homotopy quotient are different possible filling spaces. The
gluing theorem can use a chain complex of any specified choice, but
does not identify those choices or reduce their higher attachments
to a primitive slope on a single circle.

Thus the computed groups in \eqref{ntg:eq:shear-homology} are groups
of the actual nontrivial boundary mapping tori. Theorem~\ref{ntg:thm:gluing}
provides the full boundary formula for a filling once its chain
attachments are given. Neither statement supplies the unprovided
attachments of a singular fibre, nor does it imply that a resulting
six-dimensional space has sphere homology or admits a complex
structure. Those questions require the specified geometric fillings
and the comparisons in \eqref{ntg:eq:missing-equations}.
